\chapter{Mantenimiento y actualización}
\label{chap:mantenimiento}

\section{Guardar una configuración conocida}

Después de completar correctamente la primera ejecución en la tarjeta, guarda la información
básica del entorno. Servirá como referencia si una actualización provoca algún problema.

\begin{lstlisting}[language=bash]
mkdir -p evidence/known-good

akd1000-bringup doctor --json \
  > evidence/known-good/doctor.json

python -m pip freeze \
  > evidence/known-good/pip-freeze.txt

sha256sum models/*.fbz \
  > evidence/known-good/model-hashes.txt
\end{lstlisting}

Conserva también la versión del proyecto y del controlador utilizados. Si más adelante
el sistema deja de funcionar, esta configuración permite comparar qué ha cambiado.

\section{Actualizar el entorno}

Las actualizaciones del kernel, del controlador, de Python o de Akida pueden afectar al
funcionamiento de la tarjeta. Siempre que sea posible, cambia un componente cada vez y
repite después la prueba básica.

Un procedimiento sencillo es:

\begin{enumerate}
  \item comprueba la compatibilidad de la nueva versión;
  \item realiza la actualización;
  \item reinicia el equipo si es necesario;
  \item repite la detección de la tarjeta y las pruebas de software y hardware.
\end{enumerate}

Si aparece un problema, compara el sistema con la configuración guardada anteriormente
y utiliza el procedimiento del \cref{chap:diagnostico}.

El controlador utiliza DKMS para reconstruir el módulo después de una actualización del
kernel. Esto solo funciona mientras la nueva versión del kernel siga estando soportada por
el controlador.

\section{Apagado y almacenamiento}

Antes de desconectar la tarjeta, termina las ejecuciones en curso y apaga el host de forma
normal. Después, desconecta las fuentes de alimentación antes de retirar la tarjeta o el
cable FFC.

Si el montaje va a guardarse durante un periodo largo, protege las placas frente a descargas
electrostáticas y evita doblar o dejar bajo tensión el cable FFC.